\section{Fiscal Policy and the Spending Multiplier }

Before the financial crisis in 2007/08, fiscal policy was often seen as inefficient and costly. In particular, the time-lag of the implementation of the fiscal spending and the crowding-out effect on the private sector were seen as great obstacles to the use of public spending programs. Besides the theoretical economic arguments, different models have delivered contrasting results on the effects of fiscal policy. \citet{CoenenG..2010} and \citet{Afonso.2010} compare various state-of-the-art macroeconomic models and show the wide range of different frameworks in which the spending multiplier can be analyzed. The type, design, duration and funding of the fiscal stimulus are as important as the circumstances of the economic meltdown, the interaction with monetary policy and the scope for public finance. \citet{Coenen.2010} mention that an economic crisis normally has negative consequences for the current debt position of a country, as we could observe during the euro crisis, and thus the budgetary room for maneuver can be highly at risk in a situation when fiscal packages can normally be extraordinary effective.
\par
\bigskip
New evidence on recent model-based economic research provide some new understanding on whether fiscal policy can be an adequate response in a deep recession. \citet{Eggertsson.2011} suggests that stimulating short run aggregate demand, rather than the supply side of the economy, has particularly positive effects at the zero lower bound. \citet{Davig.2009} and \citet{Christiano.2011} show that an increase in government spending can outsize the effect on output in a deep recession, especially in prolonged liquidity traps. Once the nominal interest rate hits the lower bound, there can be a strong case for a fiscal stimulus on a temporary basis to jump-start the economy and push it out of an economic downward spiral.
\par
\bigskip
The key question of the paper ``Is there a fiscal free lunch in a liquidity trap?'' concentrates on the magnitude of the government spending multiplier, the distinction between marginal and average effects and the impact of the fiscal policy on the government budget.
The paper uses a simple New Keynesian model to analyze the impulse response of two different exogenous shocks. It shows that fiscal policy can play a crucial role in stabilizing the economy during an economic crisis. One important result of the paper is that the marginal multiplier is high in a liquidity trap, but decreases quickly when the size of the government spending increases. We replicate the concept of the marginal spending multiplier, which is a step function depending on the level of government spending.
We also show that under certain conditions, an expansionary fiscal policy can in fact reduce government debt and thus be self-financing, the case in which there is a fiscal free lunch. Finally, we point out the risks of fiscal austerity measures while the economy faces a liquidity trap.
